\begin{explanationbox}
	\textbf{\textcolor{CExplanation}{一句话直觉}}
	椭圆绕轴扫出的椭球体，就像被均匀拉伸的球，体积仍然带着熟悉的 \(\frac{4}{3}\pi\)，只是半径被两个半轴“分工合作”。

	\medskip
	\textbf{\textcolor{CExplanation}{核心拆解}}
	\begin{description}[style=nextline, leftmargin=2.8em, labelsep=0.5em, itemsep=0.45em]
		\item[\textbf{要点一（截面圆盘）：}] 用垂直于 \(x\) 轴的平面去截旋转体，截面是圆，半径为当时椭圆纵坐标的绝对值 \(|y|\)，因此截面面积为 \(\pi y^{2}\)。
		\item[\textbf{要点二（积分累加）：}] 将 \(x\) 从 \(-a\) 到 \(a\) 上所有薄圆盘的体积加总，得到定积分 \(V=\int_{-a}^{a} \pi y^{2}\,dx\)。
		\item[\textbf{要点三（代数化简）：}] 利用椭圆方程 \(y^{2}=b^{2}\bigl(1-\frac{x^{2}}{a^{2}}\bigr)\) 消去 \(y\)，积分后自然得到 \(V=\frac{4}{3}\pi a b^{2}\)。
	\end{description}

	\medskip
	\textbf{\textcolor{CExplanation}{几何本质（祖暅原理）}}
	从截面积分的角度看，旋转体体积等于各个位置截面面积的累积：\(V = \int_{-a}^{a} S(x)\,dx\)。这与祖暅原理的思想一致：若两个立体在同一高度处的截面积处处相等，则它们体积相等。椭圆旋转体的每个截面圆半径按椭圆规律变化，但总面积累积后仍保持与球体相同的结构系数。

	\medskip
	\textbf{\textcolor{CExplanation}{代数意义（对称简化）}}
	积分利用了区间对称性和被积函数的偶函数性质，使计算化简为 \(\pi b^{2}\cdot 2\bigl(a-\frac{a}{3}\bigr)\)，最终 \(\frac{4}{3}\) 来自 \(2(1-\frac{1}{3})\) 的乘积。

	\medskip
	\textbf{\textcolor{CExplanation}{考点价值}}
	\begin{itemize}[leftmargin=*, itemsep=0.5em, topsep=0.4em]
		\item \textbf{考法一（直接套用公式）：} 给出标准椭圆方程，直接识别 \(a,b\)，代入体积公式。
		\item \textbf{考法二（逆向求参数）：} 已知旋转椭球体积和其中一个半轴，反求另一半轴或椭圆方程。
	\end{itemize}

	\medskip
	\textbf{\textcolor{CExplanation}{顿悟点}}

	\textbf{球体积是 \(\frac{4}{3}\pi r^{3}\)，而椭球体积是 \(\frac{4}{3}\pi a b^{2}\)。为什么 \(b\) 出现平方？因为旋转垂直于 \(x\) 轴的方向是二维的，所以 \(y\) 方向的尺度对体积的影响是平方关系，而 \(x\) 方向只沿旋转轴拉伸，体积与它成正比。}
	\begin{center}
		% C025_fig07_why_b_squared_v6_stable.tex
% 用法：主文件提前加载 \usepackage{tikz,amsmath}
% 说明：这是 TikZ 片段，不包含 \documentclass、\begin{document}、\end{document}
% 核心改进：公式框和提示框都用固定坐标画，不再让 node 自动决定大框高度。

\begin{center}
	\begin{tikzpicture}[
			x=1cm,
			y=1cm,
			>=stealth,
			line cap=round,
			line join=round,
			every node/.style={font=\small},
			mainarrow/.style={->,line width=0.95pt,blue!65!black},
			axis/.style={->,line width=0.72pt,black!65}
		]

		% =========================================================
		% 固定画布
		% =========================================================
		\path[use as bounding box] (-7.25,-4.35) rectangle (7.25,4.05);

		% =========================================================
		% 标题
		% =========================================================
		\node[
			anchor=north,
			font=\bfseries\Large,
			blue!65!black
		] at (0,3.88)
		{为什么体积公式里是 \(b^2\)？};

		% =========================================================
		% 左侧：椭圆上一点
		% =========================================================
		\node[
			font=\bfseries\large,
			blue!70!black
		] at (-5.28,2.35)
		{椭圆上一点};

		\fill[blue!5] (-5.28,0.95) ellipse (1.55 and 0.72);
		\draw[blue!70!black,line width=0.95pt] (-5.28,0.95) ellipse (1.55 and 0.72);

		% 坐标轴
		\draw[axis] (-7.00,0.95) -- (-3.45,0.95)
		node[below left=-1pt] {\(x\)};
		\draw[axis] (-5.28,-0.05) -- (-5.28,1.88)
		node[left] {\(y\)};

		% 固定横坐标 x 的竖线
		\draw[dashed,black!50,line width=0.7pt] (-4.56,0.25) -- (-4.56,1.58);
		\draw[red!75!black,line width=1.25pt] (-4.56,0.95) -- (-4.56,1.50);

		\node[
			red!75!black,
			anchor=west,
			font=\normalsize
		] at (-4.24,1.36)
		{\(r(x)=|y|\)};

		\node[black!65] at (-4.56,-0.13) {\(x\)};

		\node[
			blue!65!black,
			font=\normalsize
		] at (-5.28,-0.82)
		{\(\displaystyle \frac{x^2}{a^2}+\frac{y^2}{b^2}=1\)};

		% =========================================================
		% 中间：截面圆盘
		% =========================================================
		\node[
			font=\bfseries\large,
			green!45!black
		] at (-1.55,2.35)
		{截面圆盘};

		\fill[green!6] (-1.55,0.95) circle (0.78);
		\draw[green!50!black,line width=0.95pt] (-1.55,0.95) circle (0.78);

		\fill[green!50!black] (-1.55,0.95) circle (1.35pt);
		\draw[green!55!black,line width=1pt] (-1.55,0.95) -- (-0.77,0.95);

		\node[
			green!45!black,
			anchor=south,
			font=\normalsize
		] at (-1.16,1.02)
		{\(r(x)\)};

		\node[
			green!45!black,
			font=\normalsize
		] at (-1.55,-0.82)
		{\(\displaystyle S(x)=\pi r(x)^2\)};

		% =========================================================
		% 右侧：固定尺寸公式框
		% 不再使用自动高度 node，避免不同编译环境下高度变化
		% =========================================================
		\draw[
			draw=blue!65!black,
			rounded corners=7pt,
			line width=0.95pt,
			fill=blue!2
		] (1.80,-1.55) rectangle (6.95,2.55);

		\node[
			font=\bfseries\normalsize,
			black
		] at (4.38,2.12)
		{半径先变面积};

		\node[font=\small] at (4.38,1.25)
		{\(\displaystyle r(x)=b\sqrt{1-\frac{x^2}{a^2}}\)};

		\node[font=\small] at (4.38,0.25)
		{\(\displaystyle S(x)=\pi r(x)^2\)};

		\node[font=\small] at (4.38,-0.85)
		{\(\displaystyle S(x)=\pi b^2\left(1-\frac{x^2}{a^2}\right)\)};

		% =========================================================
		% 箭头
		% =========================================================
		\draw[mainarrow] (-3.43,0.95) -- (-2.48,0.95)
		node[midway,below=7pt,font=\scriptsize,black!70]
		{绕 \(x\) 轴旋转};

		\draw[mainarrow] (-0.77,0.95) -- (1.80,0.95)
		node[midway,above=8pt,font=\scriptsize,black!70]
		{取面积函数};

		% =========================================================
		% 底部关键提示框：固定尺寸
		% =========================================================
		\draw[
			draw=red!50!black,
			rounded corners=5pt,
			line width=0.85pt,
			fill=red!3
		] (-6.30,-3.28) rectangle (6.30,-2.15);

		\node[
			font=\normalsize,
			align=center
		] at (0,-2.58)
		{\textbf{关键：}绕 \(x\) 轴旋转时，垂直于 \(x\) 轴的截面是圆；};

		\node[
			font=\normalsize,
			align=center
		] at (0,-3.02)
		{圆面积要平方半径，所以 \(y\) 方向半轴 \(b\) 在体积公式中变成 \(b^2\)。};

		% =========================================================
		% 图注
		% =========================================================
		\node[
			font=\small,
			black!80
		] at (0,-3.95)
		{图：为什么是 \(ab^2\)，不是 \(a^2b\)，也不是 \(ab\)？};

	\end{tikzpicture}
\end{center}
	\end{center}
	\medskip
	\textbf{\textcolor{CExplanation}{使用场景}}
	\begin{itemize}[leftmargin=*, itemsep=0.5em, topsep=0.4em]
		\item \textbf{场景一（求旋转体体积）：} 题目给出椭圆的标准方程，要求计算绕 \(x\) 轴旋转所得立体体积。
		\item \textbf{场景二（物理建模）：} 在物理问题中，计算椭球形容器的容积或椭球天体的大致体积。
	\end{itemize}
\end{explanationbox}
